\section*{Chapitre \ref{ch:b1:intermolecular-forces-solvents} --- Interactions intermoléculaires et solvants}

\begin{solution}{exo:b1:intermolecular-forces-solvents:1}
Argon~: London seulement. Chlorure d'hydrogène~: Keesom, Debye et
London, London dominant (\cref{exo:b1:intermolecular-forces-solvents:8})~;
le chlore est trop gros pour des \omterm{def:b1:intermolecular-forces-solvents:hbond}{liaisons hydrogène} fortes. Eau~: les
\omterm{def:b1:intermolecular-forces-solvents:hbond}{liaisons hydrogène} dominent, avec les trois termes de van der Waals.
Méthane~: London seulement. Diiode~: London seulement, intense parce que
la molécule est grosse et très polarisable.
\end{solution}

\begin{solution}{exo:b1:intermolecular-forces-solvents:2}
Entre leurs propres molécules~: \ce{CH3OH} (\ce{O-H}), \ce{CH3NH2}
(\ce{N-H}) et \ce{HF}. \ce{CH3OCH3} et \ce{CH3F} ont des \omterm{def:b1:lewis-resonance-vsepr:lewis}{doublets non
liants} sur O ou F, mais aucun hydrogène porté par ces atomes~: ils ne
peuvent qu'accepter des \omterm{def:b1:intermolecular-forces-solvents:hbond}{liaisons hydrogène}, de l'eau par exemple.
\ce{CH4} n'en donne ni n'en accepte.
\end{solution}

\begin{solution}{exo:b1:intermolecular-forces-solvents:3}
Eau et éthanol~: \omterm{def:b1:intermolecular-forces-solvents:solvent-classes}{solvants polaires} et protiques. Propanone et
acétonitrile~: polaires et aprotiques. Dichlorométhane~: peu polaire, aprotique. Hexane~:
apolaire, aprotique.
\end{solution}

\begin{solution}{exo:b1:intermolecular-forces-solvents:4}
Le diiode est apolaire~: dans l'hexane, il échange des \omterm{def:b1:intermolecular-forces-solvents:vdw}{interactions de
London} contre des \omterm{def:b1:intermolecular-forces-solvents:vdw}{interactions de London}, alors que dans l'eau il
romprait des \omterm{def:b1:intermolecular-forces-solvents:hbond}{liaisons hydrogène} sans les remplacer. Le chlorure de
sodium est ionique~: seul un solvant de forte permittivité peut séparer
ses ions, et seul un \omterm{def:b1:intermolecular-forces-solvents:solvent-classes}{solvant polaire} peut les solvater~; l'hexane
($\varepsilon_r = 1.89$) ne fait ni l'un ni l'autre.
\end{solution}

\begin{solution}{exo:b1:intermolecular-forces-solvents:5}
Le méthoxyméthane n'a pas de liaison \ce{O-H}~: il ne peut pas donner de
\omterm{def:b1:intermolecular-forces-solvents:hbond}{liaison hydrogène} et bout le plus bas. Le méthanol et l'éthanol sont tous
deux associés par \omterm{def:b1:intermolecular-forces-solvents:hbond}{liaisons hydrogène}~; l'éthanol, avec un \ce{CH2} de
plus, a des \omterm{def:b1:intermolecular-forces-solvents:vdw}{interactions de London} plus fortes. L'acide éthanoïque forme
des dimères liés par \omterm{def:b1:intermolecular-forces-solvents:hbond}{liaisons hydrogène} et c'est le plus lourd~: il bout
le plus haut.
\end{solution}

\begin{solution}{exo:b1:intermolecular-forces-solvents:6}
Dans le vide, $F = e^2/(4\pi\varepsilon_0 r^2) = (1.602 \times
10^{-19})^2/(4\pi \times 8.854 \times 10^{-12} \times (5.00 \times
10^{-10})^2) = \qty{9.23e-10}{N}$. Dans l'eau, $F/78.4 = \qty{1.18e-11}{N}$~;
dans l'hexane, $F/1.89 = \qty{4.88e-10}{N}$. L'eau affaiblit l'attraction
41 fois plus que l'hexane~: elle peut maintenir les ions séparés,
l'hexane non.
\end{solution}

\begin{solution}{exo:b1:intermolecular-forces-solvents:7}
Mélanger la propanone à l'eau rompt des \omterm{def:b1:intermolecular-forces-solvents:hbond}{liaisons hydrogène} eau--eau mais
en crée de nouvelles entre l'eau et l'oxygène de \ce{C=O}, ainsi que des
interactions dipôle--dipôle~: le bilan est favorable. L'hexane n'offre à
l'eau que de faibles \omterm{def:b1:intermolecular-forces-solvents:vdw}{interactions de London} en échange des \omterm{def:b1:intermolecular-forces-solvents:hbond}{liaisons
hydrogène} qu'elle romprait~: les deux liquides restent séparés.
\end{solution}

\begin{solution}{exo:b1:intermolecular-forces-solvents:8}
\ce{HBr} a plus d'électrons et un nuage plus gros, plus polarisable~:
son \omterm{def:b1:intermolecular-forces-solvents:vdw}{interaction de London} dépasse celle de \ce{HCl} de plus que son
\omterm{def:b1:intermolecular-forces-solvents:vdw}{interaction de Keesom} n'est inférieure. London domine.
\end{solution}

\begin{solution}{exo:b1:intermolecular-forces-solvents:9}
\chemfig{CH_3-C(=[:60]O)-[:-60]O-[:0]H} face à son partenaire, deux
liaisons $\ce{O-H}\cdots\ce{O=C}$ fermant un cycle à huit atomes. Dans le
benzène, apolaire et aprotique, rien ne concurrence ces liaisons~: l'acide
est dimérisé et se comporte comme une particule d'environ
\qty{120}{g/mol}. Dans l'eau, chaque molécule d'acide forme plutôt des
\omterm{def:b1:intermolecular-forces-solvents:hbond}{liaisons hydrogène} avec l'eau~: les dimères se dissocient.
\end{solution}

\begin{solution}{exo:b1:intermolecular-forces-solvents:10}
Le chlorure de sodium est formé d'ions~: l'eau, de forte permittivité,
les sépare (dissociation) et les hydrate (oxygène tourné vers \ce{Na+},
hydrogène vers \ce{Cl-}). Le chlorure d'hydrogène est une \omterm{def:b1:lewis-resonance-vsepr:dipole}{molécule
polaire}~: l'eau, polaire, transforme la liaison \ce{H-Cl} en paire d'ions
\ce{H3O+ Cl-} (ionisation), sépare la paire (dissociation) et hydrate les
ions. L'hexane est apolaire et de faible permittivité~: il n'ionise ni
ne dissocie, et aucun ion ne transporte le courant.
\end{solution}

\begin{solution}{exo:b1:intermolecular-forces-solvents:11}
$(174.1 - 36.1)/5 = \qty{27.6}{\celsius}$ par \ce{CH2}. Le dernier
incrément (du nonane au décane) vaut \qty{23.4}{\celsius}~: l'undécane
devrait bouillir vers $174 + 22 \approx \qty{196}{\celsius}$. Chaque
\ce{CH2} apporte la même contribution de London, mais une fraction de
plus en plus petite du total de la molécule, alors que la température
nécessaire pour vaincre les interactions est déjà élevée.
% ledger: bp:C9H20
\end{solution}

\begin{solution}{exo:b1:intermolecular-forces-solvents:12}
Partie \omterm{def:b1:intermolecular-forces-solvents:amphiphile}{hydrophile}~: la tête sulfate \ce{-OSO3^-} avec son contre-ion
\ce{Na+}~; partie \omterm{def:b1:intermolecular-forces-solvents:amphiphile}{hydrophobe}~: la chaîne à douze carbones. À la
surface, les têtes dans l'eau et les queues dans l'air~; dans une
micelle, les queues à l'intérieur et les têtes à l'extérieur. Les queues
se dissolvent dans la graisse de la tache, les têtes la maintiennent au
contact de l'eau, et la graisse est décollée dans des micelles.
\end{solution}

\begin{solution}{pb:b1:intermolecular-forces-solvents:1}
\textbf{1.} L'\omterm{def:b1:intermolecular-forces-solvents:vdw}{interaction de London}, la seule entre atomes apolaires.
\textbf{2.} De l'hélium au xénon, le nombre d'électrons et la taille
augmentent, donc la \omterm{def:b1:periodicity:polarisability}{polarisabilité} et l'attraction de London.
\textbf{3.} L'hélium a deux électrons fortement retenus par son noyau~:
c'est l'atome le moins polarisable, avec l'attraction de London la plus
faible.
\textbf{4.} $68.8 - 36.1 = \qty{32.7}{\celsius}$~: un \ce{CH2} apporte
des électrons et de la surface de contact, donc de l'attraction de
London.
\textbf{5.} London, qui croît avec la \omterm{def:b1:periodicity:polarisability}{polarisabilité}~: $\ce{CH4} < \ce{SiH4} <
\ce{GeH4}$.
\textbf{6.} Le carbone n'est pas assez électronégatif pour polariser
fortement \ce{C-H}, et le méthane n'a pas de \omterm{def:b1:lewis-resonance-vsepr:lewis}{doublet non liant} pour
accepter une \omterm{def:b1:intermolecular-forces-solvents:hbond}{liaison hydrogène}.
\textbf{7.} L'eau~: polaire, protique, de la plus forte permittivité~;
elle sépare et solvate les deux ions.
\textbf{8.} $78.4/1.89 = 41$~: l'attraction est 41 fois plus faible dans
l'eau.
\textbf{9.} L'acétonitrile (ou la propanone)~: assez polaire pour
dissoudre le sel, aprotique pour que l'anion reste libre.
\textbf{10.} L'hexane~: non \omterm{def:b1:intermolecular-forces-solvents:miscibility}{miscible} à l'eau, apolaire comme le diiode.
Le diiode passe dans l'hexane, la phase supérieure (l'hexane est moins
dense que l'eau).
\textbf{11.} L'éthanol est \omterm{def:b1:intermolecular-forces-solvents:miscibility}{miscible} à l'eau~: il ne se forme pas de
seconde phase.
\textbf{12.} Son oxygène accepte des \omterm{def:b1:intermolecular-forces-solvents:hbond}{liaisons hydrogène} de l'eau, mais il
ne peut en donner aucune et ses deux groupes éthyle sont \omterm{def:b1:intermolecular-forces-solvents:amphiphile}{hydrophobes}~: il
ne se dissout que faiblement.
\textbf{13.} $-66.4 - (-85.0) = \qty{18.6}{\celsius}$ par \omterm{def:b1:periodicity:table}{période}~;
\ce{HF} extrapolé~: $-85.0 - 18.6 = \qty{-103.6}{\celsius}$.
\textbf{14.} \ce{HF} bout à \qty{19.6}{\celsius}, \qty{123}{\celsius}
au-dessus de l'extrapolation.
\textbf{15.} $-62.5 - (-87.8) = \qty{25.3}{\celsius}$~; \ce{NH3}
extrapolé~: \qty{-113.1}{\celsius}~; valeur réelle $-33.4$, soit
\qty{80}{\celsius} de plus.
\textbf{16.} $-88.1 - (-112.0) = \qty{23.9}{\celsius}$~; \ce{CH4}
extrapolé~: \qty{-135.9}{\celsius}~; valeur réelle $-161.5$, en dessous.
Le méthane ne forme pas de \omterm{def:b1:intermolecular-forces-solvents:hbond}{liaisons hydrogène}~; il est simplement petit
et peu polarisable.
\textbf{17.} L'eau (environ \qty{179}{\celsius}) $>$ \ce{HF} (123) $>$
\ce{NH3} (80).
\textbf{18.} \ce{NH3} peut en donner trois mais n'en accepter qu'une (un
seul \omterm{def:b1:lewis-resonance-vsepr:lewis}{doublet non liant})~; \ce{HF} peut en accepter trois mais n'en donner
qu'une. L'eau en donne deux et en accepte deux~: chaque molécule peut
participer à quatre \omterm{def:b1:intermolecular-forces-solvents:hbond}{liaisons hydrogène}, un réseau tridimensionnel
complet.
\textbf{19.} \ce{H2S} et \ce{H2Se} ne forment pas de \omterm{def:b1:intermolecular-forces-solvents:hbond}{liaisons hydrogène}
notables~: leurs températures d'ébullition suivent la tendance de London
que l'eau suivrait sans elles. L'eau elle-même est l'anomalie que l'on
mesure.
\textbf{20.} $-41.3 - (-60.3) = \qty{19.0}{\celsius}$ par \omterm{def:b1:periodicity:table}{période}.
\textbf{21.} $T(p) = -60.3 + 19.0\,(p - 3)$ en \unit{\celsius}.
\textbf{22.} $T(2) = -60.3 - 19.0 = \qty{-79.3}{\celsius}$.
\textbf{23.} $100.0 - (-79.3) = \qty{179}{\celsius}$.
\textbf{24.} Sans ses \omterm{def:b1:intermolecular-forces-solvents:hbond}{liaisons hydrogène}, l'eau bouillirait vers
\textbf{\qty{-79}{\celsius}}, et serait un gaz sur Terre.
\end{solution}
